\section{Introduction}
\label{sec:introduction}

Code-switching, the alternation between two languages within a conversation or a
single sentence, is the norm rather than the exception in multilingual
communities~\cite{sitaram2019survey}. In Indian professional and academic
settings, meetings mix Hindi and English freely: a sentence such as
``\textit{release Friday ko hoga, testing kal tak}'' (``the release will be on
Friday, testing by tomorrow'') carries its decision in English nouns and its
commitment in Hindi grammar. Automatic meeting assistants, which transcribe a
meeting, attribute speech to speakers and distil it into a summary and action
items, are now common for English, but three properties of Hindi--English
meetings make them a poor fit for existing tools.

First, \emph{the script is ambiguous}. Hindi words may be written in Devanagari
or in Latin script, and a speech recognizer may emit either, even within one
utterance. A transcript that is correct word for word can still score close to
100\,\% word error rate if its reference uses the other script. Second,
\emph{downstream steps interact}. Transliterating Hindi into Latin script, which
makes a transcript readable to colleagues who do not read Devanagari, destroys
the evidence a script-based language identifier relies on; done in the wrong
order, every segment silently becomes ``English''. Third, \emph{recordings are
often confidential}, so an assistant that uploads audio to a cloud service is
frequently not an option.

This paper presents \system{}, a meeting assistant for Hindi--English
code-mixed speech that runs entirely on one machine with open-weight models, and
a reproducible evaluation of it. Our contributions are:

\begin{enumerate}
  \item \textbf{An integrated, fully local pipeline}
    (Section~\ref{sec:system}) combining Whisper large-v3~\cite{radford2023whisper},
    pyannote speaker diarization~\cite{bredin2020pyannote,plaquet2023powerset},
    script-based language identification and a locally served
    7B-parameter language model~\cite{qwen2024qwen25} to produce speaker-attributed,
    language-tagged transcripts, Romanized Hinglish, summaries and action items,
    for both recorded and live meetings.
  \item \textbf{Design rules for code-mixed pipelines}, enforced in code rather
    than by convention: language identification precedes transliteration;
    transliteration is restricted to script conversion and may not correct
    recognition errors; and the evaluation uses a deterministic transliterator,
    never the system's own, so the system is not used to score itself.
  \item \textbf{A provenance design that makes results auditable}
    (Section~\ref{sec:provenance}): each run records the code commit, a hash of
    every output-affecting setting, prompt versions and the input's checksum;
    runs whose stages were produced under different configurations are excluded
    from evaluation automatically; and all tables and figures in this paper are
    generated from stored runs by a single command.
  \item \textbf{An evaluation study} (Sections~\ref{sec:evaluation}
    and~\ref{sec:results}) covering seven families of metrics (recognition,
    diarization, cross-chunk speaker identity, language identification,
    summarization, action-item extraction and system efficiency) on
    \pending{N} annotated code-mixed meetings, with annotator agreement,
    bootstrap confidence intervals and ablations over the choices that most
    affect quality and cost.
\end{enumerate}

The code, the annotation protocol and the generated results are released at
\pending{repository URL}.
